%----------------------------------------------------------------------------------------
%    PACKAGES AND THEMES
%----------------------------------------------------------------------------------------

\documentclass[aspectratio=169,xcolor=dvipsnames]{beamer}
\usetheme{SimpleDarkBlue}

\usepackage{hyperref}
\usepackage{graphicx} % Allows including images
\usepackage{booktabs} % Allows the use of \toprule, \midrule and \bottomrule in tables
\usepackage[ruled,vlined,linesnumbered]{algorithm2e}
\usepackage{algcompatible}
\usepackage{caption}    % if you want captions inside beamer
\usepackage{makecell} % Add this to your preamble

%----------------------------------------------------------------------------------------
%    TITLE PAGE
%----------------------------------------------------------------------------------------

\title{Meta-WebGuard: An Adaptive Hybrid Framework for Cross-Domain
Web Attack Detection}

\author{\small Project Domain: Machine Learning\\
Project Guide: Dr. Paul P Mathai\\
Team Members: \\
Rahul K J (FIT22CS155)\\
Sethulakshmi P A (FIT22CS170)\\
Sidharth C J (FIT22CS176)\\
Susan Joy (FIT22CS184)\\[0.5cm]
\today
}

%----------------------------------------------------------------------------------------
%    PRESENTATION SLIDES
%----------------------------------------------------------------------------------------

\begin{document}

\begin{frame}
    \titlepage
\end{frame}

\begin{frame}{Abstract}
\begin{itemize}
    \item Implemented \textbf{Meta-WebGuard}, a multi-layered hybrid framework combining Calibrated SVM, \textbf{CNN-BiLSTM}, and an XGBoost Meta-Learner.
    \item Integrated a \textbf{Titan-98 Global Mix} dataset, bridging detection gaps using real-world "unseen" samples and CAPEC-aligned injection data.
    \item Developed a \textbf{Hybrid CNN-BiLSTM} architecture to capture both local character motifs and long-range sequential patterns.
    \item Achieved exceptional ensemble accuracy (98.81\%) and a validated \textbf{99.23\% Precision} on 100\% unseen data samples.
    \item Optimized for production using \textbf{cost-sensitive learning}, applying 2.5x weighting to malicious samples to minimize false negatives.
\end{itemize}
\end{frame}

\begin{frame}{Introduction}
\begin{itemize}
    \item \textbf{Evolving Threat Landscape:} Modern attacks use sophisticated obfuscation; our system uses \texttt{urllib} decoding to expose hidden payloads.
    \item \textbf{Hybrid Modeling:} Combines the linguistic precision of \textbf{TF-IDF SVM} with the spatial-temporal memory of \textbf{CNN-BiLSTM}.
    \item \textbf{Our Solution:} A stacked "Ensemble" approach that aggregates intelligence from statistical, textual, and neural models.
    \item \textbf{Meta-Learning:} Employs \textbf{XGBoost} to arbitrate between foundation models, identifying the most reliable indicators for each URL.
\end{itemize}
\end{frame}

\begin{frame}{Problem Statement}
\begin{itemize}
    \item \textbf{Detection Gap:} Traditional firewalls struggle with brand-spoofing; our script identifies "Look-alike" domains using a \textbf{High-Trust Whitelist}.
    \item \textbf{Static vs. Dynamic:} Fixed systems are bypassed by entropy-rich payloads; we utilize \textbf{Shannon Entropy} metrics to flag randomized strings.
    \item \textbf{Class Imbalance:} Malicious URLs are often rarer; our implementation uses \textbf{Cost-Sensitive learning} (2.5x penalty) to ensure robust detection.
    \item \textbf{Architecture Mismatch:} Standard models miss local character motifs; we solve this with a unified \textbf{CNN-BiLSTM} stack.
\end{itemize}
\end{frame}

\begin{frame}{Objectives}
\begin{itemize}
    \item Build a hybrid pipeline integrating \textbf{LinearSVC, CNN-BiLSTM, and XGBoost} for 4-class URL classification.
    \item Achieve \textbf{$>$98\% Accuracy} on the "Unseen" data subset to validate the system's Zero-Day readiness.
    \item Utilize \textbf{Calibrated Probability Scoring} to provide high-precision detection (Proven: \textbf{99.23\% Precision}).
    \item Develop a \textbf{Functional Web Interface} for real-time URL testing and visualization of detection results.
\end{itemize}
\end{frame}

\begin{frame}{Social Relevance}
\begin{itemize}
    \item Securing web applications protects users' personal information, financial data, and online identity from cybercriminals.
    \item Early detection of web attacks reduces data breach risk for individuals, businesses, and government portals.
    \item Strengthens trust and security of e-commerce platforms, healthcare portals, and social networks.
    \item Contributes to broader cybersecurity resilience by raising the bar for attack detection capability.
\end{itemize}
\end{frame}

\section*{Sustainable Development Goals (SDGs) Addressed}
\begin{frame}{Sustainable Development Goals (SDGs) Addressed}
\subsection*{SDG 9: Industry, Innovation and Infrastructure}
\begin{itemize}
    \item Strengthens cybersecurity infrastructure using AI-based attack detection.
    \item Enhances resilience of web systems against phishing and injection attacks.
\end{itemize}

\subsection*{SDG 16: Peace, Justice and Strong Institutions}
\begin{itemize}
    \item Reduces cybercrime through intelligent threat detection.
    \item Protects users and institutions from online fraud and data breaches.
\end{itemize}

\subsection*{SDG 8: Decent Work and Economic Growth}
\begin{itemize}
    \item Prevents financial losses caused by phishing and malicious websites.
    \item Enhances trust in digital commerce and online transactions.
\end{itemize}
\end{frame}

\begin{frame}{Overview of literature reviews}
\begin{table}[htbp]
\centering
\renewcommand{\arraystretch}{1.3}
\caption{Overview of Literature Review}
\resizebox{\textwidth}{!}{
\begin{tabular}{|p{1.2cm}|p{2.2cm}|p{2.5cm}|p{2.5cm}|p{2.8cm}|p{2.8cm}|p{2.8cm}|}
\hline
\textbf{Year} & \textbf{Authors} & \textbf{Paper Title} & \textbf{Objectives} & \textbf{Methodology} & \textbf{Conclusion / Results} & \textbf{Contributions / Findings} \\
\hline
2023 & Demirel et al. & Zero-Shot CNN for Web Attacks & Detect zero-day attacks & CNN-based zero-shot detection & High recall & Unseen attack pattern detection \\
\hline
2023 & Bensaid et al. & CNN-BiLSTM IDS & Improve SDN IDS accuracy & Hybrid CNN + BiLSTM & Accuracy $>$ 98\% & Spatial-temporal traffic modeling \\
\hline
2024 & Alashhab et al. & Online ML Ensemble & Real-time DDoS detection & Adaptive online ensemble & High detection rate & Dynamic adaptation capability \\
\hline
2025 & Guo et al. & Web-FTP Transfer Model & Cross-dataset generalization & Feature transfer learning & Improved accuracy & Effective with limited labels \\
\hline
\end{tabular}
}
\end{table}
\end{frame}

% Rest of the Literature tables (Truncated for brevity in summary, but full in code)
\begin{frame}{Overview of literature reviews (Cont.)}
\begin{table}[htbp]
\centering
\renewcommand{\arraystretch}{1.3}
\resizebox{\textwidth}{!}{
\begin{tabular}{|p{1.2cm}|p{2.2cm}|p{2.5cm}|p{2.5cm}|p{2.8cm}|p{2.8cm}|p{2.8cm}|}
\hline
2024 & Shayegan et al. & SDN--NFV Defense & Network DDoS prevention & SDN + NFV integration & Reduced attack impact & Flexible programmable mitigation \\
\hline
2024 & Jagat et al. & VLSTM-AE for Web Logs & Detect anomalies & Variational LSTM Autoencoder & High accuracy & Temporal deviation modeling \\
\hline
2025 & Wu et al. & Blockchain Security & Improve web resilience & Blockchain + adversarial learning & Enhanced resilience & Distributed adaptive framework \\
\hline
2025 & Kulkarni et al. & ML vs LLM Study & evaluate phishing robustness & Comparative ML vs LLM & Accuracy drops under attack & Identified robustness gaps \\
\hline
\end{tabular}
}
\end{table}
\end{frame}

\begin{frame}{System Architecture}
\begin{figure}[h!]
   \centering
  \includegraphics[width=0.67\textwidth]{system architecture.png}
  \caption{Architecture Diagram of Meta-WebGuard} \label{fig:system_architecture}
 \end{figure}
\end{frame}

\begin{frame}{Phases of Project - Neural Evolution}
\textbf{Phase 4: Hybrid CNN-BiLSTM Neural Architecture} \\[0.2cm]
Deep learning model focused on multi-scale feature extraction:
\begin{itemize}
    \item \textbf{Spatial Layer:} \textbf{Conv1D} captures local character motifs (signatures).
    \item \textbf{Temporal context:} Two stacked \textbf{Bidirectional LSTM} layers for long-range intent.
    \item \textbf{Regularization:} SpatialDropout1D and standard Dropout (0.2--0.3) for model robustness.
    \item \textbf{Ensemble Role:} Passes high-fidelity neural probabilities to the XGBoost Meta-Learner.
\end{itemize}
\end{frame}

\begin{frame}{Algorithm: Calibrated SVM \& CNN-BiLSTM (3/5)}
\begin{algorithm}[H]
\small
\begin{algorithmic}
  \State \textbf{Step 4: Train Foundation Model A (Calibrated SVM)}
  \State \hspace{1.2em}(a) Wrap LinearSVC in CalibratedClassifierCV for probability scores
  \State \hspace{1.2em}(b) Fit on character-level TF-IDF sparse matrix

  \State \textbf{Step 5: Train Foundation Model B (CNN-BiLSTM)}
  \State \hspace{1.2em}(a) Layer 1: 64-dim Character Embedding
  \State \hspace{1.2em}(b) Layer 2: \textbf{Conv1D (64 filters, kernel=3)} for spatial pattern extraction
  \State \hspace{1.2em}(c) Layer 3-4: Stacked \textbf{Bidirectional LSTM} for temporal context
\end{algorithmic}
\end{algorithm}
\end{frame}

\begin{frame}{Results: Master Verification (Unseen Data)}
\begin{table}[htbp]
\centering
\caption{Final Performance Benchmark (100\% Unseen URLs)}
\renewcommand{\arraystretch}{1.1}
\begin{tabular}{|l|c|c|}
\hline
\textbf{Metric Type} & \textbf{Achieved Value} & \textbf{Benchmark} \\
\hline
Overall Accuracy & 97.70\% & Exceeds 95\% \\
\hline
Precision (Malicious) & \textbf{99.23\%} & \textbf{Excellent} \\
\hline
Recall (Malicious) & 96.15\% & Exceeds 95\% \\
\hline
F1-Score & 97.66\% & Target: 96\% \\
\hline
\textbf{Ensemble Accuracy} & \textbf{98.81\%} & \textbf{Global Mix} \\
\hline
\end{tabular}
\end{table}
\end{frame}

\begin{frame}{Conclusion}
\begin{itemize}
    \item \textbf{Integrated Ensemble:} Successfully combined SVM (Textual), \textbf{CNN-BiLSTM} (Sequential), and XGBoost (Meta).
    \item \textbf{Verified Generalization:} Documented **99.23\% Precision** on 100\% unseen data samples.
    \item \textbf{Zero-Day Ready:} The XGBoost Meta-Learner provides an adaptive framework capable of catching novel anomalies.
    \item \textbf{Ready for Defense:} Final technical specs 100\% aligned with the production implementation.
\end{itemize}
\end{frame}

\end{document}
